% GENERATED FILE. Edit canonical lessons, then run npm run generate:books.
% canonical-source-hash: fnv1a64:ba463390537afd5e
% canonical-lessons: TA-C235-ticai, TA-C235-vatakku, TA-C235-terku, TA-C235-kilakku, TA-C235-merku

\chapter{Direction, North, South, East, West}
\label{ch:ta-direction-north-south-east-west}
\hlchaptermodality{235}

\begin{chapteropening}
\textbf{By the end of this chapter:} \emph{I can say a direction, north, south, east and west.}

Complete the last lesson of chapter 235: I can say a direction, north, south, east and west.
\end{chapteropening}

\section[\texorpdfstring{\ta{திசை} (ticai)}{திசை (ticai)}]{\ta{திசை} (ticai) — a direction}

\label{lesson:TA-C235-ticai}



\begin{quote}
Before the new one: say the Tamil for a capital city, then the Tamil for a crossroads.
\end{quote}

\subsection*{You'll want to know: \ta{திசை}}
\textbf{\ta{திசை}} — \emph{ticai} — \textquotedblleft{}a direction\textquotedblright{}.

\begin{grammarlens}[title={What you've built}]
One more word for this chapter.
\end{grammarlens}

\subsection*{Your turn}
\begin{itemize}
\raggedright
  \item \emph{Say it:} \emph{ticai}
  \item \emph{Say it:} \emph{ticai}, once more
  \item \emph{Say it:} say \emph{ticai}
  \item \emph{Recall:} say the Tamil for a capital city, then the Tamil for a crossroads, then say \emph{ticai} again
\end{itemize}

\begin{tcolorbox}[breakable,colback=teal!4,colframe=teal!35!black,title={Before you move on}]
What does \ta{திசை} mean? (\textquotedblleft{}A direction\textquotedblright{}.) Say it once more.
\end{tcolorbox}
\section[\texorpdfstring{\ta{வடக்கு} (vaṭakku)}{வடக்கு (vaṭakku)}]{\ta{வடக்கு} (vaṭakku) — north}

\label{lesson:TA-C235-vatakku}



\begin{quote}
Before the new one: say the Tamil for a crossroads, then the Tamil for a direction.
\end{quote}

\subsection*{You'll want to know: \ta{வடக்கு}}
\textbf{\ta{வடக்கு}} — \emph{vaṭakku} — \textquotedblleft{}north\textquotedblright{}.

\begin{grammarlens}[title={What you've built}]
One more word for this chapter.
\end{grammarlens}

\subsection*{Your turn}
\begin{itemize}
\raggedright
  \item \emph{Say it:} \emph{vaṭakku}
  \item \emph{Say it:} \emph{vaṭakku}, once more
  \item \emph{Say it:} say \emph{vaṭakku}
  \item \emph{Recall:} say the Tamil for a crossroads, then the Tamil for a direction, then say \emph{vaṭakku} again
\end{itemize}

\begin{tcolorbox}[breakable,colback=teal!4,colframe=teal!35!black,title={Before you move on}]
What does \ta{வடக்கு} mean? (\textquotedblleft{}North\textquotedblright{}.) Say it once more.
\end{tcolorbox}
\section[\texorpdfstring{\ta{தெற்கு} (teṟku)}{தெற்கு (teṟku)}]{\ta{தெற்கு} (teṟku) — south}

\label{lesson:TA-C235-terku}



\begin{quote}
Before the new one: say the Tamil for a direction, then the Tamil for north.
\end{quote}

\subsection*{You'll want to know: \ta{தெற்கு}}
\textbf{\ta{தெற்கு}} — \emph{teṟku} — \textquotedblleft{}south\textquotedblright{}.

\begin{grammarlens}[title={What you've built}]
One more word for this chapter.
\end{grammarlens}

\subsection*{Your turn}
\begin{itemize}
\raggedright
  \item \emph{Say it:} \emph{teṟku}
  \item \emph{Say it:} \emph{teṟku}, once more
  \item \emph{Say it:} say \emph{teṟku}
  \item \emph{Recall:} say the Tamil for a direction, then the Tamil for north, then say \emph{teṟku} again
\end{itemize}

\begin{tcolorbox}[breakable,colback=teal!4,colframe=teal!35!black,title={Before you move on}]
What does \ta{தெற்கு} mean? (\textquotedblleft{}South\textquotedblright{}.) Say it once more.
\end{tcolorbox}
\section[\texorpdfstring{\ta{கிழக்கு} (kiḻakku)}{கிழக்கு (kiḻakku)}]{\ta{கிழக்கு} (kiḻakku) — east}

\label{lesson:TA-C235-kilakku}



\begin{quote}
Before the new one: say the Tamil for north, then the Tamil for south.
\end{quote}

\subsection*{You'll want to know: \ta{கிழக்கு}}
\textbf{\ta{கிழக்கு}} — \emph{kiḻakku} — \textquotedblleft{}east\textquotedblright{}.

\begin{grammarlens}[title={What you've built}]
One more word for this chapter.
\end{grammarlens}

\subsection*{Your turn}
\begin{itemize}
\raggedright
  \item \emph{Say it:} \emph{kiḻakku}
  \item \emph{Say it:} \emph{kiḻakku}, once more
  \item \emph{Say it:} say \emph{kiḻakku}
  \item \emph{Recall:} say the Tamil for north, then the Tamil for south, then say \emph{kiḻakku} again
\end{itemize}

\begin{tcolorbox}[breakable,colback=teal!4,colframe=teal!35!black,title={Before you move on}]
What does \ta{கிழக்கு} mean? (\textquotedblleft{}East\textquotedblright{}.) Say it once more.
\end{tcolorbox}
\section[\texorpdfstring{\ta{மேற்கு} (mēṟku)}{மேற்கு (mēṟku)}]{\ta{மேற்கு} (mēṟku) — west}

\label{lesson:TA-C235-merku}



\begin{quote}
Before the new one: say the Tamil for south, then the Tamil for east.
\end{quote}

\subsection*{You'll want to know: \ta{மேற்கு}}
\textbf{\ta{மேற்கு}} — \emph{mēṟku} — \textquotedblleft{}west\textquotedblright{}.

\begin{grammarlens}[title={What you've built}]
One more word for this chapter.
\end{grammarlens}

\subsection*{Your turn}
\begin{itemize}
\raggedright
  \item \emph{Say it:} \emph{mēṟku}
  \item \emph{Say it:} \emph{mēṟku}, once more
  \item \emph{Say it:} say \emph{mēṟku}
  \item \emph{Recall:} say the Tamil for south, then the Tamil for east, then say \emph{mēṟku} again
\end{itemize}

\begin{tcolorbox}[breakable,colback=teal!4,colframe=teal!35!black,title={Before you move on}]
What does \ta{மேற்கு} mean? (\textquotedblleft{}West\textquotedblright{}.) Say it once more.
\end{tcolorbox}
